\chapter*{Abstract}
\addcontentsline{toc}{chapter}{Abstract}
Sequential quasi-Monte Carlo (SQMC) is an extension of sequential Monte Carlo (SMC) in which independent Monte Carlo inputs are replaced by randomised low-discrepancy point sets and particles are ordered before resampling; under regularity conditions it can attain lower integration error than SMC, up to the Monte Carlo rate of $O_P(N^{-1/2})$. However, SQMC is more computationally expensive than SMC at $\mathcal{O}(N \log N)$ and its advantage may deteriorate as the effective dimension grows. We develop an implementation of SQMC on a GPU using \texttt{jax}, and measure its accelerator performance: on a linear-Gaussian filtering model an A100 GPU executed SQMC 6.2--13.4 times faster than a CPU at 2,048 particles across the five dimensions. We then demonstrate how Rao-Blackwellisation can be used to tackle high-dimensional limitations of SQMC in the context of football matches. This particle approach does not, however, outperform a factorial state-space model with an extended Kalman filter (EKF) approximation, indicating that the EKF is a suitable approximation for this problem. This work is part of a contribution to the open-source state-space inference package \texttt{cuthbert}.

\textbf{Keywords}: Sequential Monte Carlo, Sequential quasi-Monte Carlo, Low-discrepancy point sets, Hilbert sort, Rao-Blackwellisation, Parallel accelerators, Football